\section{Method}

Figure~\ref{fig:method} summarizes the post-hoc verifier. Given a GUI screenshot and target cue, SeekUI first predicts a scanpath and target status. We then compute cognitive path evidence from whether the scanpath visits target-like candidates, and OCR candidate evidence from fuzzy matching between the query and visible text. A combined AND rule rejects a target only when both evidence sources are weak. We also evaluate an evidence-aware VLM that receives the screenshot and structured path/OCR evidence.

\begin{figure}[t]
  \centering
  \includegraphics[width=\linewidth]{method_diagram.png}
  \caption{Post-hoc target-presence verification for GUI visual search. The verifier converts forced-choice grounding into uncertainty-aware present/absent decisions without retraining SeekUI.}
  \label{fig:method}
\end{figure}

\subsection{Prompt-Only Baseline}

The first baseline uses SeekUI's prompt-only output. If the model emits an absent status, we evaluate it as absent; otherwise, we treat the prediction as present. This baseline represents the behavior of a target-present visual-search model when it is asked to handle target absence without additional evidence checks.

\subsection{Cognitive Path Evidence}

The cognitive stopping layer measures whether the predicted scanpath provides evidence for a target-like candidate. Intuitively, a present target should eventually attract the path toward a stable region, while an absent target often yields weak, diffuse, or unsupported evidence. The path score is not a new scanpath generator; it is a post-hoc signal derived from the generated path.

For each screenshot, we build a candidate set from known target annotations in the released data. Each candidate has a center point $c_j$ and text label $t_j$. This candidate inventory is annotation-backed: it is appropriate for controlled benchmark analysis because it lets us ask whether the generated scanpath visits any known target-like UI element, but it is not itself a deployable perception module. A deployed version would replace this inventory with OCR boxes, accessibility-tree elements, UI detector proposals, icon proposals, or crop-level VLM candidates. For a query $q$ and predicted fixation point $f_i$, we compute:

\[
  e(f_i, c_j, q, t_j) =
  \operatorname{sim}(q,t_j) \cdot \exp\left(-\frac{\|f_i-c_j\|_2}{r}\right),
\]

where $\operatorname{sim}$ is a normalized fuzzy string similarity and $r$ is a radius proportional to the image diagonal. The path evidence score is the maximum evidence over all predicted fixations and candidates:

\[
  E_{\mathrm{path}} = \max_{i,j} e(f_i,c_j,q,t_j).
\]

A low $E_{\mathrm{path}}$ increases the likelihood that the model is making a forced-choice commitment rather than finding a real target.

\subsection{OCR Candidate Verification}

Because the released subset is text-target dominated, OCR provides an interpretable UI evidence source. We extract OCR candidates from the screenshot and compare their text to the target cue using fuzzy matching. OCR evidence is useful when the target is textual and visibly rendered, but it is brittle for small, stylized, low-contrast, cropped, or non-text targets.

The OCR verifier therefore should not be treated as a complete GUI understanding model. Its role is narrower: it checks whether visible text supports the target cue.

Let $\mathcal{O}(I)$ be OCR candidates extracted from screenshot $I$ after filtering by OCR confidence. The OCR evidence score is:

\[
  E_{\mathrm{ocr}} = \max_{o_k \in \mathcal{O}(I)} \operatorname{sim}(q,o_k).
\]

\subsection{Combined AND Verifier}

The combined verifier rejects a target only when both path evidence and OCR evidence are weak:

\[
  \hat{y} =
  \begin{cases}
  \mathrm{absent}, & E_{\mathrm{path}} < \tau_{\mathrm{path}} \land E_{\mathrm{ocr}} < \tau_{\mathrm{ocr}}\\
  \mathrm{present}, & \mathrm{otherwise.}
  \end{cases}
\]

This AND rule is intentionally conservative: it tries to reduce absent false-present errors while avoiding unnecessary rejection of visible targets. Thresholds are selected on development examples by optimizing absent F1, then evaluated on held-out examples. The image-disjoint split is the primary evaluation because it avoids shared screenshots across development and test sets.

\begin{figure}[t]
\small
\fbox{\begin{minipage}{0.94\linewidth}
\textbf{Algorithm 1: Combined path+OCR verifier}

\begin{enumerate}
  \item Run SeekUI on screenshot $I$ and query $q$ to obtain scanpath $S$ and original status.
  \item Build benchmark UI candidates from target annotations and OCR candidates from the screenshot.
  \item Compute $E_{\mathrm{path}} = \max_{f_i,c_j} \operatorname{sim}(q,t_j)\exp(-\|f_i-c_j\|_2/r)$.
  \item Compute $E_{\mathrm{ocr}} = \max_{o_k} \operatorname{sim}(q,o_k)$ over OCR candidates above the confidence threshold.
  \item If both $E_{\mathrm{path}} < \tau_{\mathrm{path}}$ and $E_{\mathrm{ocr}} < \tau_{\mathrm{ocr}}$, output \texttt{absent}; otherwise output \texttt{present}.
\end{enumerate}
\end{minipage}}
\caption{Reproducible decision rule for the interpretable combined verifier.}
\label{fig:combined_algorithm}
\end{figure}

\subsection{VLM Baselines and Evidence-Aware VLM}

We also evaluate generic VLM yes/no target-presence prompts. These baselines answer a reviewer-risk question: could we simply ask a VLM whether the target is visible? We test direct, conservative, OCR-aware, and search-behavior prompts.

Finally, we evaluate an evidence-aware VLM verifier. This verifier sees the screenshot as primary evidence and also receives structured path/OCR evidence from the post-hoc pipeline. We report it both on the full benchmark and on the matched image split. In the current framing, combined AND is the main interpretable verifier and evidence-aware VLM is the strongest validated verifier.

The prompt is intentionally constrained so the result is not an unconstrained chat completion. A compact excerpt is:

\begin{quote}
\small
You are checking whether the target is visible in the UI screenshot. Use the screenshot as primary evidence. Scanpath/OCR evidence is diagnostic but may be wrong. Answer only: \texttt{status: present} or \texttt{status: absent}.
\end{quote}

\begin{figure}[t]
\small
\fbox{\begin{minipage}{0.94\linewidth}
\textbf{Algorithm 2: Evidence-aware VLM verifier}

\begin{enumerate}
  \item Start from the combined-verifier record for screenshot $I$ and query $q$.
  \item Construct a VLM prompt containing the screenshot, target cue, original status, combined status, $E_{\mathrm{path}}$, $E_{\mathrm{ocr}}$, best OCR text, and thresholds.
  \item Instruct the VLM to use the screenshot as primary evidence and answer only \texttt{status: present} or \texttt{status: absent}.
  \item Parse the answer using explicit present/absent patterns; if parsing is ambiguous, default to \texttt{present}.
\end{enumerate}
\end{minipage}}
\caption{Evidence-aware VLM verifier. The prompt includes structured evidence but asks the VLM to judge the screenshot as the primary source.}
\label{fig:vlm_algorithm}
\end{figure}
